\FloatBarrier
\section*{Выводы по главе 3}
\phantomsection\addcontentsline{toc}{section}{Выводы по главе 3}

В главе~3 описана численная реализация решения задач гидродинамики
смазочного слоя, включая дискретизацию уравнения Рейнольдса,
итерационные методы, GPU-ускорение и реализацию модели кавитации JFO.

\begin{enumerate}[wide, labelindent=\parindent, nosep]
  \item Сформулированы граничные условия для дискретной задачи:
    условия Дирихле ($p = p_{\mathrm{cav}}$) на открытых границах,
    периодичность по окружному направлению (в частности, через ghost-узлы
    в GPU-реализации), условия Неймана при расчёте сектора.
    В реализации граничные условия выделены в отдельный этап.
  \item Выполнена консервативная дискретизация уравнения Рейнольдса
    на основе баланса потоков через грани контрольного объёма.
    Напорные потоки аппроксимированы центральными разностями
    (второй порядок), а переносная часть в постановке JFO
    (перенос $\vartheta$) реализована upwind-схемой первого порядка,
    обеспечивающей устойчивость при доминировании конвективного переноса.
    Дискретизация для давления (при $\vartheta=1$) приводит к линейной
    системе с пятиточечным шаблоном.
  \item Для решения дискретной системы выбран метод последовательной
    верхней релаксации (SOR) как естественное ускорение метода
    Гаусса--Зейделя.  Сформулированы критерии сходимости
    (относительная поправка и невязка разностного уравнения)
    и определены параметры итерационного процесса.
  \item Для преодоления ограничения последовательного SOR
    на параллелизм реализован метод Red-Black SOR на GPU.
    Шахматная раскраска сетки и двухполуходная схема (red/black)
    обеспечивают независимость обновлений внутри полухода, что позволяет
    эффективно использовать массово-параллельную архитектуру GPU.
    Описаны конфигурация запуска ядер, кэширование GPU-буферов для серийных
    расчётов и механизм контроля сходимости с минимальными накладными
    расходами на глобальную редукцию.
  \item Реализован алгоритм решения задачи кавитации JFO
    с двухуровневой итерационной структурой: внутренний цикл
    (GPU Red-Black SOR для давления при фиксированном~$\vartheta$)
    и внешний цикл (обновление активного множества и пересчёт степени
    заполнения~$\vartheta$ из баланса переносных потоков).
    Обновление~$\vartheta$ в кавитационной зоне основано на дискретном
    уравнении переноса куэттовским потоком с upwind-аппроксимацией.
    Реализация при консервативной аппроксимации потоков обеспечивает
    массосохраняющее решение задачи кавитации.
\end{enumerate}

Разработанный вычислительный алгоритм (дискретизация, GPU-решатель, JFO)
применяется ко всем рассматриваемым постановкам (упорный, радиальный
и возвратно-поступательный подшипники) и используется в расчётах,
представленных в последующих главах.
